\section{Perlatan dan Kakas}

Peralatan dan kakas yang digunakan dalam pengembangan aplikasi Tamu.id dan proyek akhir ini adalah sebagai berikut:

\subsection{Perangkat Lunak}

\begin{enumerate}
    \item Visual Studio Code: Editor kode sumber yang digunakan untuk menulis kode program.
    \item Flask: Framework Python yang digunakan untuk membangun backend aplikasi.
    \item VueJS:Library JavaScript yang digunakan untuk membangun frontend aplikasi.
    \item Tailwind CSS: Framework CSS yang digunakan untuk mendesain antarmuka pengguna.
    \item PrimeVue: Komponen UI yang digunakan untuk mempercepat pengembangan antarmuka pengguna terutama bagian Super Admin dan Admin.
    \item Axios: Library JavaScript yang digunakan untuk melakukan permintaan HTTP dari frontend ke backend.
    \item JSON Web Token (JWT): Standar terbuka yang digunakan untuk mengamankan komunikasi antara frontend dan backend.
    \item Docker: Platform yang digunakan untuk mengemas aplikasi dan dependensinya ke dalam kontainer, sehingga memudahkan pengembangan, pengujian, dan penyebaran.
    \item GitLab: Platform pengembangan perangkat lunak yang digunakan untuk mengelola kode sumber.
    \item PostgreSQL: Sistem manajemen basis data relasional yang digunakan untuk menyimpan data aplikasi.
    \item Google Cloud Run: serverless platform tempat deploy aplikasi backend.
    \item Google Cloud Storage: Layanan penyimpanan objek yang digunakan untuk menyimpan file statis.
    \item Google Cloud SQL: Layanan basis data terkelola yang digunakan untuk menyimpan data aplikasi di lingkungan Produksi.
    \item Cloudflare Pages: Layanan hosting statis yang digunakan untuk menyajikan aplikasi frontend.
    \item Midtrans Core API: Layanan pembayaran yang digunakan untuk memproses transaksi pembayaran.
    \item Ngrok: Alat yang digunakan untuk membuat terowongan ke localhost, memungkinkan akses ke aplikasi yang sedang dikembangkan di mesin lokal.
    \item Hoppscotch: Alat pengujian API yang digunakan untuk menguji endpoint API.
    \item PostgreSQL: Sistem manajemen basis data relasional yang digunakan untuk menyimpan data aplikasi.
    \item K6: Platform pengujian beban yang digunakan untuk menguji performa aplikasi.
    \item Latex: Bahasa markup yang digunakan untuk menulis dokumen proyek akhir.
    \item Draw.io (VSCode Extension): Ekstensi Visual Studio Code yang digunakan untuk membuat diagram alur.
    \item PlantUML: Alat yang digunakan untuk membuat diagram UML, seperti diagram kelas, diagram urutan, dan diagram aktivitas.
    \item Mermaid: Alat yang digunakan untuk membuat diagram alur, diagram kelas, dan diagram lainnya dalam format teks.
    \item  ClickUp: Alat manajemen proyek yang digunakan untuk mengelola tugas-tugas dalam proyek akhir ini.
\end{enumerate}

\subsection{Perangkat Keras}

Berikut spesifikasi perangkat keras yang digunakan dalam pengembangan aplikasi AtleTiket dan proyek akhir ini:

\begin{enumerate}
    \item Laptop Thinkpad T480s
          \begin{itemize}
              \item Prosesor: Intel Core i5-8350U
              \item RAM: 16 GB DDR4 2400 MHz
              \item Penyimpanan: 256 GB SSD NVMe M.2
              \item Kartu Grafis: Intel UHD Graphics
              \item Sistem Operasi: Ubuntu 24.02 LTS
          \end{itemize}
    \item Smartphone Infinix Hot 9
          \begin{itemize}
              \item Prosesor: MediaTek Helio A25
              \item RAM: 4 GB
              \item Penyimpanan: 128 GB
              \item Sistem Operasi: Android 10
              \item Browser: Google Chrome
              \item Resolusi Layar: 720 x 1600 piksel
              \item Ukuran Layar: 6.6 inci
          \end{itemize}
\end{enumerate}